\begin{assumption}[Smoothness and strong convexity]\label{ass:curv}
The fine-tuning objective $F:\mathbb{R}^d\to\mathbb{R}$ is $L$-smooth and $\mu$-strongly convex on $\mathbb{R}^d$, with $0<\mu\le L<\infty$. It has a unique minimizer $x^\star:=\arg\min_x F(x)$, $F^\star:=F(x^\star)$, and $\nabla F(x^\star)=0$. Equivalently, the Bregman divergence $D_F(x,y):=F(x)-F(y)-\langle\nabla F(y),x-y\rangle$ obeys $\tfrac{\mu}{2}\|x-y\|^2\le D_F(x,y)\le\tfrac{L}{2}\|x-y\|^2$ for all $x,y$.
\end{assumption}

\begin{definition}[DP-SGD update]\label{def:dpsgd}
Let $N$ be the dataset size and $\mathcal D=\{\ell_1,\dots,\ell_N\}$ the per-example losses. Fix a clip norm $C>0$, expected batch size $B$ (Poisson rate $q=B/N$), and noise multiplier $\sigma_{\mathrm{DP}}\ge0$. Writing $\mathrm{clip}_C(v):=v\cdot\min\{1,C/\|v\|_2\}$, the privatized gradient at iterate $x$ is
\[
g_{\mathrm{DP}}(x)\;=\;\frac1B\Big[\sum_{i\in\mathcal B}\mathrm{clip}_C\big(\nabla\ell_i(x)\big)\;+\;z\Big],
\qquad z\sim\mathcal N\!\big(0,(\sigma_{\mathrm{DP}}C)^2 I_d\big),
\]
where $\mathcal B$ is a Poisson($q$) subsample and $z$ is independent of $\mathcal B$ and of the data. DP-SGD runs $x^{k+1}=x^k-\gamma\,g_{\mathrm{DP}}(x^k)$ with constant step $\gamma>0$. The per-coordinate DP-noise variance baked into $g_{\mathrm{DP}}$ is
\[
\boxed{\;\Phi\;:=\;\Big(\frac{\sigma_{\mathrm{DP}}\,C}{B}\Big)^2\;}
\qquad\Longrightarrow\qquad \mathbb E\big\|z/B\big\|^2=d\,\Phi,
\]
which matches the optimizer code exactly (\texttt{phi = (noise\_multiplier*max\_grad\_norm/B\_eff)**2}, with $B_{\mathrm{eff}}=$ expected batch $\times$ accumulation $\times$ world size).
\end{definition}

\begin{assumption}[Unbiasedness / small-clip regime]\label{ass:unbiased}
On the trajectory the per-sample gradients satisfy $\|\nabla\ell_i(x)\|\le C$, so clipping acts as the identity, and Poisson subsampling is unbiased. Hence $\mathbb E[g_{\mathrm{DP}}(x)\mid x]=\nabla F(x)$. (When clipping bites, an additive bias $b(x)$ appears; this is treated in the clipping-bias remark and only inflates the floor, so the statement below is a lower bound on the true floor.)
\end{assumption}

\begin{assumption}[Expected smoothness / ABC inequality]\label{ass:es}
There is a finite constant $A>0$ (the \emph{expected-smoothness} constant) and the optimum-variance $\sigma_\star^2:=\mathbb E\|g_{\mathrm{DP}}(x^\star)\|^2$ such that, for all $x$,
\[
\mathbb E\big\|g_{\mathrm{DP}}(x)-g_{\mathrm{DP}}(x^\star)\big\|^2\;\le\;2A\,\big(F(x)-F^\star\big).
\]
This is the standard unified-SGD / expected-smoothness hypothesis (Gower et al.\ 2019; Khaled \& Richt\'arik 2020). The optimum variance splits, by independence of the Gaussian noise from the clipped signal, as
\[
\sigma_\star^2\;=\;\sigma_{\mathrm{sub}}^2\;+\;d\,\Phi,
\qquad\text{where }\sigma_{\mathrm{sub}}^2:=\mathbb E\|\bar g_{\mathrm{clip}}(x^\star)\|^2\text{ is the subsampling variance.}
\]
\end{assumption}

\begin{theorem}[DP-SGD converges to a noise floor with an irreducible $d\Phi$ DP term]\label{thm:floor}
Under Assumptions~\ref{ass:curv}--\ref{ass:es}, DP-SGD with any constant step
\[
0<\gamma\le\frac{1}{2A}
\]
satisfies, for every $k\ge0$,
\[
\mathbb E\|x^k-x^\star\|^2\;\le\;(1-\gamma\mu)^k\,\|x^0-x^\star\|^2\;+\;\frac{2\gamma\,\sigma_\star^2}{\mu}\,.
\]
Consequently $\limsup_{k\to\infty}\mathbb E\|x^k-x^\star\|^2\le \dfrac{2\gamma\,\sigma_\star^2}{\mu}$, and substituting the variance split of Assumption~\ref{ass:es} isolates the differential-privacy contribution to the floor,
\[
\boxed{\;\limsup_{k\to\infty}\mathbb E\|x^k-x^\star\|^2
\;\le\;\underbrace{\frac{2\gamma\,\sigma_{\mathrm{sub}}^2}{\mu}}_{\text{subsampling}}
\;+\;\underbrace{\frac{2\gamma\,d\,\Phi}{\mu}}_{=:\ \mathrm{FLOOR}_{\mathrm{DP}}}\,,
\qquad
\mathrm{FLOOR}_{\mathrm{DP}}=\frac{2\gamma\,d\,\sigma_{\mathrm{DP}}^2 C^2}{\mu\,B^2}\;}.
\]
The DP floor is linear in the step $\gamma$, linear in the trainable dimension $d$ (so DP-LoRA, which replaces $d$ by the LoRA parameter count, lowers it), grows with the noise multiplier $\sigma_{\mathrm{DP}}$ (stricter $\varepsilon$) and clip $C$, shrinks like $1/B^2$ in the batch, and is irreducible: it is independent of $x$ and does not vanish at $x^\star$, so no number of iterations can anneal it.
\end{theorem}

\begin{proof}
We prove Theorem~\ref{thm:floor} in four steps: (i) a second-moment (``$AC$'') bound on $g_{\mathrm{DP}}$; (ii) a one-step contraction of $\mathbb E\|x^k-x^\star\|^2$; (iii) unrolling the recursion to the stated bound; (iv) substituting the DP-variance split to expose $\mathrm{FLOOR}_{\mathrm{DP}}\propto d\Phi$. Throughout, $\mathbb E_k[\,\cdot\,]:=\mathbb E[\,\cdot\mid x^k]$ denotes expectation over the fresh Poisson subsample and Gaussian noise of step $k$, conditioned on $x^k$; $\mathbb E[\cdot]$ is the total expectation. All steps use only Assumptions~\ref{ass:curv}--\ref{ass:es}.

\medskip
\noindent\textbf{Preliminary fact ($A\ge\mu$).}
We record one inequality used in Step~2 to control the contraction factor. Apply Assumption~\ref{ass:es} to the special case where the stochastic gradient is the full (noise-free, full-batch) gradient $g=\nabla F$, for which $g(x^\star)=\nabla F(x^\star)=0$; this gives $\|\nabla F(x)\|^2\le 2A\,(F(x)-F^\star)$. On the other hand, $\mu$-strong convexity implies the Polyak--\L ojasiewicz inequality $\|\nabla F(x)\|^2\ge 2\mu\,(F(x)-F^\star)$ for all $x$ (apply \eqref{eq:scvx} below at $x$ with Cauchy--Schwarz and $\|\nabla F(x)\|\,\|x-x^\star\|\ge\langle\nabla F(x),x-x^\star\rangle\ge (F(x)-F^\star)+\tfrac\mu2\|x-x^\star\|^2\ge 2\sqrt{\tfrac\mu2(F-F^\star)}\,\|x-x^\star\|\cdot\tfrac1{\sqrt2}$, the standard derivation). Chaining,
\[
2\mu\,(F(x)-F^\star)\;\le\;\|\nabla F(x)\|^2\;\le\;2A\,(F(x)-F^\star)
\qquad\Longrightarrow\qquad A\ge\mu,
\]
taking any $x$ with $F(x)>F^\star$. Hence the prescribed step obeys $\gamma\le\tfrac1{2A}\le\tfrac1{2\mu}$, so $\gamma\mu\le\tfrac12$ and the contraction factor $r:=1-\gamma\mu\in[\tfrac12,1)$ is in particular nonnegative and strictly less than $1$. (If one prefers to take Assumption~\ref{ass:es} as a black box, note that the unified-SGD literature always has $A\ge\mu$; either way $r\in[\tfrac12,1)$.)

\medskip
\noindent\textbf{Step 1 (Second-moment / $AC$ bound).}
We first bound $\mathbb E\|g_{\mathrm{DP}}(x)\|^2$ by the suboptimality gap. For any vectors $a,b$, the parallelogram-type inequality $\|a+b\|^2\le 2\|a\|^2+2\|b\|^2$ (from $\|a-b\|^2\ge0$) gives, with $a=g_{\mathrm{DP}}(x)-g_{\mathrm{DP}}(x^\star)$ and $b=g_{\mathrm{DP}}(x^\star)$ (so $a+b=g_{\mathrm{DP}}(x)$, and the inequality holds pointwise for each shared realization of subsample and noise),
\[
\big\|g_{\mathrm{DP}}(x)\big\|^2\;\le\;2\big\|g_{\mathrm{DP}}(x)-g_{\mathrm{DP}}(x^\star)\big\|^2\;+\;2\big\|g_{\mathrm{DP}}(x^\star)\big\|^2.
\]
Taking expectations and applying the expected-smoothness inequality (Assumption~\ref{ass:es}) to the first term and the definition $\sigma_\star^2=\mathbb E\|g_{\mathrm{DP}}(x^\star)\|^2$ to the second,
\begin{equation}\label{eq:AC}
\mathbb E\big\|g_{\mathrm{DP}}(x)\big\|^2\;\le\;\underbrace{2(2A)}_{=:\,\mathcal A}\big(F(x)-F^\star\big)\;+\;\underbrace{2\sigma_\star^2}_{=:\,\mathcal C}\,,
\qquad \mathcal A=4A,\ \ \mathcal C=2\sigma_\star^2.
\end{equation}
This is the standard $\mathbb E\|g\|^2\le \mathcal A\,(F-F^\star)+\mathcal C$ inequality: the stochastic-gradient second moment is controlled by the \emph{suboptimality gap} plus an additive constant $\mathcal C$ that does \emph{not} vanish as $x\to x^\star$.

\medskip
\noindent\textbf{Step 2 (One-step contraction).}
Fix $k$ and expand the squared distance after the update $x^{k+1}=x^k-\gamma\,g_{\mathrm{DP}}(x^k)$:
\[
\|x^{k+1}-x^\star\|^2
=\|x^k-x^\star\|^2-2\gamma\big\langle g_{\mathrm{DP}}(x^k),\,x^k-x^\star\big\rangle+\gamma^2\big\|g_{\mathrm{DP}}(x^k)\big\|^2.
\]
Apply $\mathbb E_k[\cdot]$. By unbiasedness (Assumption~\ref{ass:unbiased}), $\mathbb E_k[g_{\mathrm{DP}}(x^k)]=\nabla F(x^k)$, and since $x^k-x^\star$ is $x^k$-measurable, $\mathbb E_k\langle g_{\mathrm{DP}}(x^k),x^k-x^\star\rangle=\langle\nabla F(x^k),x^k-x^\star\rangle$. Hence
\begin{equation}\label{eq:onestep0}
\mathbb E_k\|x^{k+1}-x^\star\|^2
=\|x^k-x^\star\|^2-2\gamma\big\langle\nabla F(x^k),x^k-x^\star\big\rangle+\gamma^2\,\mathbb E_k\big\|g_{\mathrm{DP}}(x^k)\big\|^2.
\end{equation}
\emph{Lower bound on the inner product.} By $\mu$-strong convexity, for all $x$,
\[
F(x^\star)\ge F(x)+\langle\nabla F(x),x^\star-x\rangle+\tfrac{\mu}{2}\|x^\star-x\|^2,
\]
which rearranges (negate the inner product, move $F(x)-F(x^\star)$ across) to the standard ``three-point'' inequality
\begin{equation}\label{eq:scvx}
\big\langle\nabla F(x),x-x^\star\big\rangle\;\ge\;\big(F(x)-F^\star\big)+\tfrac{\mu}{2}\|x-x^\star\|^2 .
\end{equation}
\emph{Upper bound on the second moment.} By \eqref{eq:AC} with $x=x^k$,
\[
\mathbb E_k\big\|g_{\mathrm{DP}}(x^k)\big\|^2\le 4A\big(F(x^k)-F^\star\big)+2\sigma_\star^2 .
\]
Substituting both bounds into \eqref{eq:onestep0},
\begin{align*}
\mathbb E_k\|x^{k+1}-x^\star\|^2
&\le \|x^k-x^\star\|^2-2\gamma\Big[\big(F(x^k)-F^\star\big)+\tfrac{\mu}{2}\|x^k-x^\star\|^2\Big]
+\gamma^2\Big[4A\big(F(x^k)-F^\star\big)+2\sigma_\star^2\Big]\\[2pt]
&=(1-\gamma\mu)\|x^k-x^\star\|^2
-\underbrace{\big(2\gamma-4A\gamma^2\big)}_{=\,2\gamma(1-2A\gamma)}\big(F(x^k)-F^\star\big)
+2\gamma^2\sigma_\star^2 .
\end{align*}
The step-size cap $\gamma\le 1/(2A)$ gives $1-2A\gamma\ge0$, and $F(x^k)-F^\star\ge0$, so the middle term is $\le0$ and may be dropped. This yields the contraction
\begin{equation}\label{eq:contract}
\mathbb E_k\|x^{k+1}-x^\star\|^2\;\le\;(1-\gamma\mu)\,\|x^k-x^\star\|^2\;+\;2\gamma^2\sigma_\star^2 .
\end{equation}
By the Preliminary fact the contraction factor satisfies $r:=1-\gamma\mu\in[\tfrac12,1)$. Taking total expectation $\mathbb E[\cdot]$ of \eqref{eq:contract} via the tower property,
\begin{equation}\label{eq:recursion}
\mathbb E\|x^{k+1}-x^\star\|^2\;\le\;(1-\gamma\mu)\,\mathbb E\|x^k-x^\star\|^2\;+\;2\gamma^2\sigma_\star^2 .
\end{equation}

\medskip
\noindent\textbf{Step 3 (Unrolling the recursion).}
Write $r:=1-\gamma\mu$, $a_k:=\mathbb E\|x^k-x^\star\|^2$, and $c:=2\gamma^2\sigma_\star^2\ge0$, so \eqref{eq:recursion} reads $a_{k+1}\le r\,a_k+c$, with $r\in[\tfrac12,1)$ established above. We claim by induction that
\begin{equation}\label{eq:claim}
a_k\;\le\;r^k a_0+c\sum_{j=0}^{k-1}r^j,\qquad k\ge0,
\end{equation}
with the empty sum at $k=0$ equal to $0$. The base case $a_0\le a_0$ is trivial. Assuming \eqref{eq:claim} for $k$, and using $r\ge0$ so that multiplication by $r$ preserves the inequality,
\[
a_{k+1}\le r\,a_k+c\le r\Big(r^k a_0+c\sum_{j=0}^{k-1}r^j\Big)+c
=r^{k+1}a_0+c\sum_{j=0}^{k-1}r^{j+1}+c
=r^{k+1}a_0+c\sum_{j=0}^{k}r^{j},
\]
which is \eqref{eq:claim} for $k+1$. Since $0\le r<1$, every term $r^j\ge0$, so the partial geometric sum is bounded by its (convergent) limit,
\[
\sum_{j=0}^{k-1}r^j\;\le\;\sum_{j=0}^{\infty}r^j\;=\;\frac{1}{1-r}\;=\;\frac{1}{\gamma\mu}
\]
(this monotone upper bound is exactly where $r\ge0$, i.e.\ $\gamma\le1/(2\mu)$, is needed; for $r<0$ the partial sums would not be monotone in $k$). Plugging this and $r=1-\gamma\mu$, $c=2\gamma^2\sigma_\star^2$ into \eqref{eq:claim},
\begin{equation}\label{eq:final}
\mathbb E\|x^k-x^\star\|^2\;\le\;(1-\gamma\mu)^k\,\|x^0-x^\star\|^2\;+\;\frac{2\gamma^2\sigma_\star^2}{\gamma\mu}
\;=\;(1-\gamma\mu)^k\,\|x^0-x^\star\|^2\;+\;\frac{2\gamma\,\sigma_\star^2}{\mu}\,,
\end{equation}
which is the first displayed bound of the theorem (using $a_0=\|x^0-x^\star\|^2$, a constant). Letting $k\to\infty$, the geometric term $(1-\gamma\mu)^k\to0$ (as $0\le r<1$), leaving
\begin{equation}\label{eq:limsup}
\limsup_{k\to\infty}\mathbb E\|x^k-x^\star\|^2\;\le\;\frac{2\gamma\,\sigma_\star^2}{\mu}\,.
\end{equation}
Thus DP-SGD contracts geometrically at rate $(1-\gamma\mu)$ toward $x^\star$, but only into a ball of squared radius $R^2=2\gamma\sigma_\star^2/\mu$ --- the \emph{variance error floor}. The step-size cap $\gamma\le 1/(2A)$ is exactly what makes the expected-smoothness term in Step~2 nonpositive and hence absorbable, and (via $A\ge\mu$) simultaneously guarantees $r\ge0$.

\medskip
\noindent\textbf{Step 4 (Isolating the $d\Phi$ DP term).}
It remains to split the floor by the structure of $\sigma_\star^2$. Evaluate $g_{\mathrm{DP}}$ at $x^\star$ and decompose it into its clipped-mean part and its Gaussian part:
\[
g_{\mathrm{DP}}(x^\star)=\bar g_{\mathrm{clip}}(x^\star)+\frac{z}{B},\qquad
\bar g_{\mathrm{clip}}(x^\star):=\frac1B\sum_{i\in\mathcal B}\mathrm{clip}_C\big(\nabla\ell_i(x^\star)\big),\quad z\sim\mathcal N(0,(\sigma_{\mathrm{DP}}C)^2 I_d).
\]
Because $z$ is drawn independently of the subsample $\mathcal B$ and of the data, and $\mathbb E[z]=0$, the cross term vanishes:
\[
\mathbb E\big\langle \bar g_{\mathrm{clip}}(x^\star),\,z/B\big\rangle
=\big\langle \mathbb E[\bar g_{\mathrm{clip}}(x^\star)],\,\mathbb E[z/B]\big\rangle=0 .
\]
Hence the second moment is exactly additive:
\[
\sigma_\star^2=\mathbb E\|g_{\mathrm{DP}}(x^\star)\|^2
=\underbrace{\mathbb E\|\bar g_{\mathrm{clip}}(x^\star)\|^2}_{=:\,\sigma_{\mathrm{sub}}^2}
+\underbrace{\mathbb E\|z/B\|^2}_{=\,d\,\Phi}.
\]
For the Gaussian term, $\mathbb E\|z/B\|^2=\tfrac1{B^2}\sum_{j=1}^d \mathbb E[z_j^2]=\tfrac1{B^2}\cdot d\,(\sigma_{\mathrm{DP}}C)^2=d\,\big(\sigma_{\mathrm{DP}}C/B\big)^2=d\,\Phi$, matching Definition~\ref{def:dpsgd}. Substituting $\sigma_\star^2=\sigma_{\mathrm{sub}}^2+d\Phi$ into \eqref{eq:limsup},
\[
\limsup_{k\to\infty}\mathbb E\|x^k-x^\star\|^2
\;\le\;\frac{2\gamma}{\mu}\big(\sigma_{\mathrm{sub}}^2+d\,\Phi\big)
\;=\;\frac{2\gamma\,\sigma_{\mathrm{sub}}^2}{\mu}+\underbrace{\frac{2\gamma\,d\,\Phi}{\mu}}_{\mathrm{FLOOR}_{\mathrm{DP}}},
\]
and expanding $\Phi=(\sigma_{\mathrm{DP}}C/B)^2$ gives $\mathrm{FLOOR}_{\mathrm{DP}}=2\gamma\,d\,\sigma_{\mathrm{DP}}^2C^2/(\mu B^2)$, as claimed. The term $d\Phi$ is constant in $x$ (it is the variance the Gaussian mechanism injects regardless of where the iterate is) and strictly positive whenever $\sigma_{\mathrm{DP}}>0$, so it persists even though $\nabla F(x^\star)=0$: it is an \emph{irreducible} contribution that no amount of iteration can remove. This is the convergence-theory shadow of the classical strongly-convex DP-ERM rate $\widetilde O\!\big(d/(N^2\varepsilon^2)\big)$, whose $d$ and $1/\varepsilon^2$ dependence enter entirely through this additive DP-noise variance, and which is matched by information-theoretic lower bounds (Bassily, Smith \& Thakurta 2014). \hfill$\blacksquare$

\medskip
\noindent\textbf{Coupling to the privacy budget (corollary).}
To express the floor in $\varepsilon$, use the moments/RDP accountant: achieving $(\varepsilon,\delta)$-DP over $T$ steps at rate $q=B/N$ requires $\sigma_{\mathrm{DP}}\ge c\,q\sqrt{T\log(1/\delta)}/\varepsilon$ for an absolute constant $c$ (Abadi et al.\ 2016; the RDP/PRV accountants of Mironov 2017 used in our implementation give the same scaling with a smaller constant). Substituting, and using $q=B/N$ so that $q^2/B^2=1/N^2$,
\[
\Phi=\Big(\frac{\sigma_{\mathrm{DP}}C}{B}\Big)^2
=\frac{c^2 (B/N)^2 T\log(1/\delta)}{\varepsilon^2}\cdot\frac{C^2}{B^2}
=\frac{c^2 C^2 T\log(1/\delta)}{N^2\varepsilon^2},
\]
so the explicit batch dependence \emph{cancels} and
\[
\mathrm{FLOOR}_{\mathrm{DP}}(\varepsilon)=\frac{2\,c^2\,\gamma\,d\,C^2\,T\log(1/\delta)}{\mu\,N^2\,\varepsilon^2}\;\propto\;\varepsilon^{-2},
\]
recovering the monotone $\varepsilon^{-2}$ utility degradation (and $\mathrm{FLOOR}_{\mathrm{DP}}\to0$ as $\varepsilon\to\infty$, where it collapses to the clip-only floor $2\gamma\sigma_{\mathrm{sub}}^2/\mu$). The residual dependence on $B$ enters only through $T$: at fixed epochs $E$ one has $T=E N/B$, so $\mathrm{FLOOR}_{\mathrm{DP}}\propto 1/B$ --- larger batches still lower the floor, consistent with the empirical $\rho$-saturation analysis.

\medskip
\noindent\textbf{Remark (clipping bias).}
Assumption~\ref{ass:unbiased} (clipping inactive) is what makes $g_{\mathrm{DP}}$ unbiased. When per-sample gradients exceed $C$, clipping introduces a bias $b(x):=\mathbb E[\bar g_{\mathrm{clip}}(x)]-\nabla F(x)\ne0$, which adds a term $-2\gamma\langle b(x^k),x^k-x^\star\rangle$ to \eqref{eq:onestep0}. Carrying it through with Young's inequality yields $\mathrm{FLOOR}_{\mathrm{true}}=\tfrac{2\gamma}{\mu}(\sigma_{\mathrm{sub}}^2+d\Phi)+O(\|b\|^2/\mu^2)$; since this only adds a nonnegative term, the floor of Theorem~\ref{thm:floor} is a \emph{lower bound} on the true DP-SGD floor. The bias vanishes as $C\to\infty$ or when the per-sample gradient noise is symmetric (Chen, Wu \& Hong 2020).

\medskip
\noindent\textbf{Remark (non-convex analogue).}
Under $L$-smoothness alone (no strong convexity), the same $AC$ template $\mathbb E\|g_{\mathrm{DP}}\|^2\le \mathcal A(F-F^\star)+\mathcal C$ with $\mathcal C=\mathcal C_{\mathrm{base}}+d\Phi$ gives, by the descent lemma, a stationarity bound $\tfrac1K\sum_{k<K}\mathbb E\|\nabla F(x^k)\|^2\le O\!\big(\Delta_F/(\gamma K)\big)+O\!\big(\gamma L\,\mathcal C\big)$ with $\Delta_F:=F(x^0)-F^\star$, whose stationary floor is again $\propto \gamma\,d\,\sigma_{\mathrm{DP}}^2C^2/B^2$. Thus the $d\Phi$ scaling survives the non-convex setting relevant to LLM fine-tuning; only the exponent on $\varepsilon$ changes (general-convex $\sqrt d/\varepsilon$ vs.\ strongly-convex $d/\varepsilon^2$).

\medskip
\noindent\textbf{Connection to bias correction (why the floor is the right object).}
The \emph{same} scalar $\Phi$ that appears here as the irreducible variance floor is the bias that DP-AdamBC subtracts from Adam's second moment: $\mathbb E[\hat v_t]=\hat v_t^{\mathrm{true}}+\Phi$. Bias correction removes $\Phi$ from the \emph{preconditioner geometry} but does not touch $\sigma_\star^2$, and therefore by Theorem~\ref{thm:floor} \emph{cannot} shrink $\mathrm{FLOOR}_{\mathrm{DP}}$. Reducing the floor requires smaller $\sigma_{\mathrm{DP}}$ (looser $\varepsilon$), larger $B$, or smaller $d$ (LoRA) --- exactly the levers our experiments confirm, and the reason the noise share $\rho=\Phi/\hat v$ saturating at $1$ leaves bias correction with no utility to recover.
\end{proof}
